\subsection{A brief summary of the approach}

We can summarize our approach into the following steps.

\begin{enumerate}[itemsep=1pt, topsep=0pt]
\item \textbf{Select Training and Testing Data.} Identify distinct samples of human-generated data that are presumably generated by the same mechanisms as the novel target setting(s) of interest. 
When possible, use multiple distinct datasets for both training and testing to increase confidence and minimize the possibility of selecting spurious \persona{s}.
\item \textbf{Propose Theory-Driven Candidate \persona{s}.} Generate a broad set of \persona{s} that are plausibly consistent with the proposed theory (or causal mechanisms if known) related to the training, testing, and novel settings.
\item \textbf{Optimize \persona{s} on Training Data.} Optimize the candidate \persona{s} to best match the training data. 
This might involve selecting a mixture of \persona{s} or adjusting trait parameters to minimize some statistical distance from observed responses.
Confirm that the optimized sample outperforms the baseline LLM off-the-shelf on the training data.
\item \textbf{Validate \persona{s} on Testing Data.} Apply the optimized \persona{s} to the testing data and evaluate their performance relative to the baseline LLM.
\end{enumerate}

Broadly speaking, our approach provides a disciplined ``trial run'' to confirm whether a given sample of \aisub{s} can reliably predict human behavior across multiple related settings.
Similar to applying economic models estimated from past data to inform predictions in new but structurally similar environments, the success of our approach depends critically on identifying and leveraging underlying stable behavioral relationships.

In Section~\ref{sec:predict_new_AR}, we demonstrate the approach empirically on a set of strategic games using training data from \cite{1120Arad2012}.
Appendix~\ref{app:predict_new_CR} provides another example using training data from \citep{charness2002understanding} and allocation games.
Both applications show that the approach substantially reduces prediction error relative to baseline LLM predictions in preregistered experiments with novel human samples.
We emphasize throughout that the two key elements of the approach---grounding candidate \persona{s} in economic or behavioral theories, and validating across distinct but related datasets---each independently address critical pitfalls.
Without theoretical grounding, optimized \persona{s} may fail to meaningfully improve even in-sample predictions; without validating across multiple related settings, optimized \persona{s} are prone to overfitting a single training context.
